\documentclass[12pt]{article}
\usepackage{amsmath,amssymb}

\begin{document}
\section*{{2025 AMC 12B Problems/Problem 22}}
Source: AoPS Wiki

\subsection*{Solution}
Because we can factor $z$ from each vertex, the area of the triangle is equal to $|z|^2$ multiplied by the area of the triangle with vertices $2, 1+i, 1-i$ . The vertical side of that triangle has length $2$ and the altitude to that side has length $1$ , so the area is $1$ . We now just need the maximum possible value of $|z|^2$ . The equation $|4z-2| = 1$ can be rewritten as $|z-1/2| = 1/4$ . Therefore, it exists as a circle with radius $1/4$ and center at $(1/2, 0)$ on the complex plane. The maximum possible magnitude of a complex number on this circle occurs with $\frac{3}{4}$ , with magnitude $\frac{3}{4}$ . We can justify this because the line connecting the origin and the complex number $\frac{3}{4}$ goes through the center of the circle. Our answer is then $\left(\frac{3}{4}\right)^2 = \boxed{\frac{9}{16}}.$ ~lprado

\end{document}
